\documentclass[10pt,a4paper]{amsart}
\usepackage[margin=19mm]{geometry}
\usepackage{amsmath,amssymb,mathtools}
\usepackage[scr=rsfs,cal=euler]{mathalfa}
\usepackage{xfrac}
\usepackage[unicode,hidelinks,bookmarks=false]{hyperref}
\newcommand{\ie}{i.e.\ }
\newcommand{\cf}{cf.\ }
\newcommand{\resp}{respectively\ }
\newcommand{\Subsection}{\subsection}
\providecommand{\nobreakdash}{\nobreak\hbox{-}}
\setlength{\emergencystretch}{2em}
\multlinegap0pt
\usepackage{bm}
\usepackage{bbm}
\usepackage{dsfont}
\usepackage{xspace}
\usepackage{cancel}
\usepackage{mathtools}
\usepackage{amsthm}
\makeatletter
\@ifundefined{theorem}{\newtheorem{theorem}{Theorem}}{}
\@ifundefined{case}{\newtheorem{case}{Case}}{}
\@ifundefined{lemma}{\newtheorem{lemma}[theorem]{Lemma}}{}
\makeatother
\title[Graphs missing a connected partition]{Graphs missing a connected partition}
\author{Foster Tom}
\date{Source: arXiv:2409.12934v2 (CC BY 4.0)}
\begin{document}
\maketitle

\noindent\emph{Source and licence.} Graphs missing a connected partition. Version arXiv:2409.12934v2 (CC BY 4.0), distributed under the Creative Commons Attribution 4.0 International licence, as recorded in the arXiv licence metadata for this exact version and shown on the abstract page https://arxiv.org/abs/2409.12934v2. Full rights notice: licenses/arxiv-2409.12934-CC-BY-4.0.txt.\par\smallskip

\noindent\textit{Editorial scope.} Counted unit: the Lemma whose source label is \texttt{lem:analysis} in the article (source0.tex lines 457--465, source label \texttt{lem:analysis}), delivered together with the article's own complete original proof of it (source0.tex lines 467--576, one \texttt{proof} environment of 110 lines). No mathematics is written, repaired or re-derived: statement and proof are the article's own text line for line. Required context retained uncounted: none. Every definition, display and label of those lines is retained unchanged. Omitted from this package and not redistributed: 1--456, 466--466, 577--781. Nothing inside a retained excerpt is omitted or altered: every retained body is delivered exactly as the article's lines are written, so no omission receipt of any kind accompanies this package. Citations: the retained excerpts carry no citation command at all, so no bibliography entry of the article is transcribed and no reference is redistributed. Reference adaptation: none was needed and none was made; every reference inside the delivered unit resolves inside it, so no unresolved reference or citation is produced. Printed slips kept literal: no printed word, symbol, hypothesis or number of the counted statement or of its proof was altered. Layout and class: the article's own class is \texttt{amsart} with packages including inputenc, amsmath, amsfonts, amssymb, stmaryrd, latexsym, graphicx, subfigure, color, hyperref, verbatim, xy, graphics, pdfsync; the wrapper is the lane's pinned \texttt{amsart} editorial wrapper at 10pt on A4, which restates the theorem-like environments the retained text needs and adds the article's own macro definitions that the retained text needs (none), with extra wrapper packages (bm, bbm, dsfont, xspace, cancel, mathtools). The delivered numbering is the unit's own. Assets: no retained span contains a figure environment, a graphics-inclusion command, a PDF-page inclusion, a table or a TikZ picture; no image file is copied into this package, the package declares no asset dependency and no image file is redistributed with it. Licence: version arXiv:2409.12934v2 of this article is distributed under the Creative Commons Attribution 4.0 International licence, as recorded in the arXiv licence metadata for this exact version and shown on the abstract page https://arxiv.org/abs/2409.12934v2. A full rights notice is delivered at licenses/arxiv-2409.12934-CC-BY-4.0.txt. Characters: every retained excerpt is pure ASCII, so the delivered engine, XeLaTeX with its default Latin Modern text font, reports no missing character.
\begin{lemma}\label{lem:analysis}
Fix positive integers $b$ and $c$ with $c\geq 500$ and $2c\leq b\leq\frac{c^2}2$. There exists a positive integer $q$ such that, letting
\begin{equation}
x=\left\lceil\frac{b+1}q\right\rceil\text{ and }y=\left\lfloor\frac{b+c}q\right\rfloor,
\end{equation} we have $x\geq c+1$, $x<y$, and the inequality
\begin{equation}
\left\lceil\frac{x-1}{y-x}\right\rceil x\leq 2b+c+1.
\end{equation}
\end{lemma}

\begin{proof}
We first note that
\begin{equation*}
y-x=\left\lfloor\frac{b+c}q\right\rfloor-\left\lceil\frac{b+1}q\right\rceil\geq\left(\frac{b+c}q-\frac{q-1}q\right)-\left(\frac{b+1}q+\frac{q-1}q\right)=\frac cq-2+\frac 1q\geq \frac cq-2
\end{equation*}
because the floors and ceilings round by at most $\frac{q-1}q$. We also have
\begin{equation*}
\left\lceil\frac{x-1}{y-x}\right\rceil x=\left\lceil\frac{\left\lceil\frac{b+1}q\right\rceil-1}{y-x}\right\rceil\left\lceil\frac{b+1}q\right\rceil\leq\left(\frac {b/q}{y-x}+1\right)\left(\frac bq+1\right).
\end{equation*}
Now we will have a few different cases, depending on the value of $b$.

\begin{case} $2c\leq b\leq \frac{c^2}{20}$. We let $q=\left\lfloor\frac bc\right\rfloor$, which means that $\frac bc-1\leq q\leq \frac bc$. Then
\begin{equation*}
x=\left\lceil\frac{b+1}q\right\rceil\geq\left\lceil\frac{b+1}bc\right\rceil\geq c+1.
\end{equation*}
Because $b\geq 2c$, we have $q\geq 2$ and therefore
\begin{equation*}
q\geq\frac 23(q+1)\geq \frac 23\frac bc.
\end{equation*} 
Because $b\leq \frac{c^2}{20}$, we have $q\leq\frac bc\leq\frac c{20}$, so $y-x\geq\frac cq-2\geq 18$ and 
\begin{equation*}
y-x\geq\frac {18}{20}(y-x+2)\geq\frac{18}{20}\frac cq.
\end{equation*}
Now using that $c\geq 500$, we have
\begin{align*}
\left\lceil\frac{x-1}{y-x}\right\rceil x&\leq\left(\frac {b/q}{y-x}+1\right)\left(\frac bq+1\right)\leq\left(\frac{20}{18}\frac bc+1\right)\left(\frac 32c+1\right)\leq(1.12\frac bc+1)(1.51c)\\&\leq 1.7b+0.51c+c\leq 1.7b+0.255b+c=1.955b+c\leq 2b+c+1.
\end{align*}
\end{case}

\begin{case} $\frac{c^2}{20}\leq b\leq\frac{c^2}{4.2}$. We let $q=\left\lfloor\frac bc\right\rfloor$ again, but our estimates will look a bit different. We still have $x\geq c+1$ as before. Because $b\geq\frac{c^2}{20}\geq 25c$, we have $q\geq 25$ and therefore
\begin{equation*}
q\geq \frac {25}{26}(q+1)\geq\frac{25}{26}\frac bc.
\end{equation*}
Because $b\leq\frac{c^2}{4.2}$, we have $q\leq\frac bc\leq\frac c{4.2}$. Now since $y-x$ is an integer and $y-x\geq\frac cq-2\geq 2.2$, we must have $y-x\geq 3$ and
\begin{equation*}
y-x\geq\frac 35(y-x+2)\geq \frac 35\frac cq.
\end{equation*}
Now using that $c\geq 500$, we have
\begin{align*}
\left\lceil\frac{x-1}{y-x}\right\rceil x&\leq\left(\frac{b/q}{y-x}+1\right)\left(\frac bq+1\right)\leq\left(\frac 53\frac bc+1\right)\left(\frac{26}{25}c+1\right)\leq(1.67\frac bc+1)(1.05c)\\&\leq 1.76b+0.05c+c\leq 1.76b+0.002b+c=1.762b+c\leq 2b+c+1.
\end{align*}
\end{case}

\begin{case}$\frac{c^2}{4.2}\leq b\leq\frac{c^2}{3.4}$. We let $q=\lfloor 0.46\sqrt b\rfloor$. Note that $0.487\leq \frac 1{\sqrt{4.2}}\leq\frac{\sqrt b}c\leq\frac 1{\sqrt{3.4}}\leq 0.543$. This means that
\begin{equation*}
x=\left\lceil\frac{b+1}q\right\rceil\geq\frac bq\geq \frac b{0.46\sqrt b}>\frac{\sqrt b}{0.487}\geq c,
\end{equation*}
so $x\geq c+1$ because $x$ is an integer. We also have
\begin{equation*}
c\geq\sqrt{3.4}\sqrt b\geq\frac{\sqrt{3.4}}{0.46}q\geq 4.008q,
\end{equation*}
so because $y-x$ is an integer and $y-x\geq\frac cq-2>2.008$, we must have $y-x\geq 3$ and
\begin{equation*}
y-x\geq\frac 35(y-x+2)\geq\frac 35\frac cq.
\end{equation*}
Now using that $\sqrt b\geq 0.487c\geq 0.487\cdot 500\geq 240$, we have $q\geq 0.46\sqrt b-1\geq 0.4558\sqrt b$ and
\begin{align*}
\left\lceil\frac{x-1}{y-x}\right\rceil x&\leq\left(\frac53\frac bc+1\right)\left(\frac bq+1\right)\leq\left(\frac 5{3\sqrt{3.4}}\sqrt b+1\right)\left(\frac{\sqrt b}{0.4558}+1\right)\\&\leq(0.904\sqrt b+1)(2.194\sqrt b+1)\leq (0.9082\sqrt b)(2.1982\sqrt b)\leq 1.997b\leq 2b+c+1.
\end{align*}
\end{case}

\begin{case}$\frac{c^2}{3.4}\leq b\leq\frac{c^2}2$. Let $q_0=\left\lfloor\frac{\sqrt b}{\sqrt{3.5}}\right\rfloor$. We will either take $q=q_0$ or $q=q_0-1$. We first show that for one of these values of $q$, we will have $y-x\geq 2$. Note that $\frac 1{\sqrt{3.4}}\leq\frac{\sqrt b}c\leq\frac 1{\sqrt 2}$, so
\begin{equation*}
c\geq\sqrt 2\sqrt b\geq\sqrt 2\sqrt{3.5}q_0\geq 2.64q_0.
\end{equation*} Also note that because $c\geq 500$, we have $\sqrt b\geq\frac{500}{\sqrt{3.4}}\geq 270$ and $q_0\geq\frac{270}{\sqrt{3.5}}-1\geq 140$. Let us first consider $q=q_0$. Use integer division to write $b+1=c_0q_0+r_0$ for some integers $c_0$ and $0\leq r_0\leq q_0-1$. If 
\begin{equation*}r_0\geq 0.38q_0\geq 0.37q_0+1,\end{equation*}
this will mean that $c+r_0-1\geq 2.64q_0+0.37q_0=3.01q_0>3q_0$. Therefore,
\begin{equation*}
b+c=b+1+c-1=c_0q_0+(c+r_0-1)>(c_0+3)q_0,
\end{equation*}
so at least three multiples of $q_0$, namely $(c_0+1)q_0$, $(c_0+2)q_0$, and $(c_0+3)q_0$, lie strictly between $b+1$ and $b+c$. This means that $x\leq c_0+1$ and $y\geq c_0+3$, so $y-x\geq 2$.\\

Now if $r_0<0.38q_0$, we will consider $q=q_0-1$. We have $\sqrt{3.5}q_0\leq\sqrt b\leq\sqrt{3.5}(q_0+1)$, so
\begin{equation*}
3.5q_0^2\leq b\leq 3.5q_0^2+7q_0+3.5.
\end{equation*}
Now the integer $c_0$ satisfies $b+1-q_0\leq c_0q_0\leq b+1$, so using that $q_0\geq 140$, we have
\begin{equation*}
3.49q_0^2\leq 3.5q_0^2-q_0\leq b+1-q_0\leq c_0q_0\leq b+1\leq 3.5q_0^2+7q_0+4.5\leq 3.56q_0^2.
\end{equation*}
This means that in the integer division of $c_0$ by $q_0$, we have
\begin{equation*}
c_0=3q_0+r_1,\text{ where }0.49q_0\leq r_1\leq 0.56q_0.
\end{equation*}
This allows us to calculate the integer division of $b+1$ by $q=q_0-1$ as
\begin{equation*}
b+1=c_0q_0+r_0=c_0(q_0-1)+3q_0+r_0+r_1=(c_0+3)(q_0-1)+(r_0+r_1+3),
\end{equation*}
where the remainder $r=r_0+r_1+3$ satisfies
\begin{equation*}
0.49(q_0-1)\leq r_1\leq r\leq 0.38q_0+0.56q_0+3=0.94q_0+3\leq 0.97(q_0-1).
\end{equation*}
This means that $c+r-1\geq 2.64q_0+0.49(q_0-1)-1>3(q_0-1)$. Therefore,
\begin{equation*}
b+c=b+1+c-1=(c_0+3)(q_0-1)+c+r-1>(c_0+6)(q_0-1),
\end{equation*}
so at least three multiples of $(q_0-1)$, namely $(c_0+4)(q_0-1)$, $(c_0+5)(q_0-1)$, and $(c_0+6)(q_0-1)$, lie strictly between $b+1$ and $b+c$. This means that $x\leq c_0+4$ and $y\geq c_0+6$, so $y-x\geq 2$. This proves our claim that one of $q=q_0$ or $q=q_0-1$ will have $y-x\geq 2$.\\

Now taking either $q=q_0$ or $q=q_0-1$ so that $y-x\geq 2$, we have
\begin{equation*}
x=\left\lceil\frac{b+1}q\right\rceil\geq\frac b{q_0}\geq\sqrt{3.5}\sqrt b>\sqrt{3.4}\sqrt b\geq c,
\end{equation*}
so $x\geq c+1$ because $x$ is an integer. Because $\sqrt b\geq 270$, we have 
\begin{equation*}q\geq q_0-1\geq \frac{\sqrt b}{\sqrt{3.5}}-2\geq\frac{\sqrt b}{\sqrt{3.8}}.\end{equation*} Finally, we have
\begin{align*}
\left\lceil\frac{x-1}{y-x}\right\rceil x&\leq\left(\frac b{2q}+1\right)\left(\frac bq+1\right)\leq \left(\frac {\sqrt{3.8}b}{2\sqrt b}+1\right)\left(\frac{\sqrt{3.8}b}{\sqrt b}+1\right)\\&\leq(0.975\sqrt b+1)(1.95\sqrt b+1)\leq (0.979\sqrt b)(1.954\sqrt b)\leq 1.92b\leq 2b+c+1.
\end{align*}
\end{case}
This completes the proof.
\end{proof}

\end{document}
